\documentclass[aspectratio=169,10pt]{ctexbeamer}
% 格点 QCD 千页课程统一样式。由所有分卷与全集共同输入。
\usepackage{amsmath,amssymb,mathtools,bm}
\usepackage{booktabs,tabularx,array}
\usepackage{longtable}
\usepackage{tikz}
\usetikzlibrary{arrows.meta,positioning,calc,fit,shapes.geometric}
\usepackage{pgfplots}
\pgfplotsset{compat=1.17}
\usepackage{graphicx}
\usepackage{adjustbox}
\usepackage{listings}
\usepackage{xcolor}
\usepackage{microtype}
\usepackage{hyperref}

\definecolor{courseblue}{HTML}{17365D}
\definecolor{coursecyan}{HTML}{147D92}
\definecolor{coursegreen}{HTML}{1E7A46}
\definecolor{courseorange}{HTML}{B85C00}
\definecolor{coursered}{HTML}{A62C2B}
\definecolor{coursepurple}{HTML}{6A3D9A}
\definecolor{coursemuted}{HTML}{5B6573}
\definecolor{coursepanel}{HTML}{F3F6F9}
\definecolor{courseline}{HTML}{C9D3DF}

\setbeamersize{text margin left=0.56cm,text margin right=0.56cm}
\setlength{\emergencystretch}{2em}
\setbeamertemplate{navigation symbols}{}
\setbeamertemplate{blocks}[rounded][shadow=false]
\setbeamercolor{normal text}{fg=black!86,bg=white}
\setbeamercolor{structure}{fg=courseblue}
\setbeamercolor{frametitle}{fg=courseblue,bg=white}
\setbeamercolor{block title}{fg=courseblue,bg=coursepanel}
\setbeamercolor{block body}{fg=black!86,bg=coursepanel}
\setbeamercolor{alerted text}{fg=coursered}
\setbeamerfont{frametitle}{size=\large,series=\bfseries}
\setbeamerfont{normal text}{size=\normalsize}
\setbeamertemplate{frametitle}{\vskip0.10cm\insertframetitle\par\vskip0.06cm}
\setbeamertemplate{footline}{%
  \hbox{\begin{beamercolorbox}[wd=\paperwidth,ht=2.3ex,dp=0.9ex,
    leftskip=0.56cm,rightskip=0.56cm]{author in head/foot}%
    \scriptsize\color{coursemuted}\insertshortauthor\hfill
    \insertshorttitle\hfill\insertframenumber/\inserttotalframenumber
  \end{beamercolorbox}}}

\hypersetup{unicode=true,colorlinks=true,linkcolor=courseblue,urlcolor=coursecyan}

\newcommand{\CourseAnchor}[2]{\hypertarget{#1:#2}{}}
\newcommand{\CourseID}[2]{\textcolor{courseblue}{\bfseries #1~#2}}
\newcommand{\CourseRef}[2]{\hyperlink{#1:#2}{\textcolor{coursecyan}{#1~#2}}}
\newcommand{\FrameAnchor}[1]{\CourseAnchor{Fr}{#1}}
\newcommand{\KnowledgeID}[1]{\CourseAnchor{K}{#1}\CourseID{K}{#1}}
\newcommand{\DefinitionID}[1]{\CourseAnchor{Def}{#1}\CourseID{Def}{#1}}
\newcommand{\EquationID}[1]{\CourseAnchor{Eq}{#1}\CourseID{Eq}{#1}}
\newcommand{\DerivationID}[1]{\CourseAnchor{Der}{#1}\CourseID{Der}{#1}}
\newcommand{\TheoremID}[1]{\CourseAnchor{Thm}{#1}\CourseID{Thm}{#1}}
\newcommand{\FigureID}[1]{\CourseAnchor{Fig}{#1}\CourseID{Fig}{#1}}
\newcommand{\TableID}[1]{\CourseAnchor{Tbl}{#1}\CourseID{Tbl}{#1}}
\newcommand{\AlgorithmID}[1]{\CourseAnchor{Alg}{#1}\CourseID{Alg}{#1}}
\newcommand{\ExerciseID}[1]{\CourseAnchor{Ex}{#1}\CourseID{Ex}{#1}}
\newcommand{\SolutionID}[1]{\CourseAnchor{Sol}{#1}\CourseID{Sol}{#1}}
\newcommand{\SourceID}[1]{\CourseAnchor{Src}{#1}\CourseID{Src}{#1}}
\newcommand{\SympyID}[1]{\CourseAnchor{SYM}{#1}\CourseID{SYM}{#1}}

\newcommand{\StatusImplemented}{\textcolor{coursemuted}{接口存在}}
\newcommand{\StatusTested}{\textcolor{coursecyan}{测试通过}}
\newcommand{\StatusClosed}{\textcolor{coursegreen}{方案闭合}}
\newcommand{\StatusPhysical}{\textcolor{coursepurple}{真实数据验证}}
\newcommand{\StatusUnverified}{\textcolor{courseorange}{未验证}}

\newcommand{\SourceLine}[1]{%
  \vspace{0.04em}\par{\raggedright\scriptsize\color{coursemuted}来源：#1\par}}

\newcommand{\BoundaryLine}[1]{%
  \vspace{0.04em}\par{\raggedright\scriptsize\color{courseorange}适用边界：#1\par}}

\newcommand{\PhysicalFlow}[4]{%
  \begin{center}
  \begin{tikzpicture}[node distance=0.42cm and 0.42cm]
    \tikzset{courseflow/.style={draw=coursecyan,fill=coursecyan!7,
      rounded corners=2pt,align=center,text width=3.15cm,minimum height=0.84cm,
      font=\small,inner sep=3pt},
      coursearrow/.style={-{Latex[length=2mm]},draw=courseblue,semithick}}
    \node[courseflow] (a) {#1};
    \node[courseflow,right=of a] (b) {#2};
    \node[courseflow,right=of b] (c) {#3};
    \draw[coursearrow] (a) -- (b);
    \draw[coursearrow] (b) -- (c);
  \end{tikzpicture}
  \par\vspace{0.08cm}{\scriptsize\color{coursemuted}\FigureID{#4}：概念关系示意；颜色只辅助分层。}
  \end{center}}

\newcommand{\CheckMark}{\textcolor{coursegreen}{\bfseries\ensuremath{\checkmark}}}
\newcommand{\RiskMark}{\textcolor{courseorange}{\bfseries !}}
\newcommand{\CodePath}[1]{\texttt{\detokenize{#1}}}
\newsavebox{\CourseEquationMeasureBox}
\newcommand{\CourseDisplayEquation}[1]{%
  \sbox{\CourseEquationMeasureBox}{$\displaystyle #1$}%
  \ifdim\wd\CourseEquationMeasureBox>1.30\linewidth
    \PackageError{lqcd-course}{Main equation would be scaled below the readability floor}%
      {Split the equation with aligned/gathered at a mathematical boundary; do not hide it with tiny text.}%
  \fi
  \begin{adjustbox}{max width=0.96\linewidth,max totalheight=0.16\textheight}%
    $\displaystyle #1$%
  \end{adjustbox}}

\lstset{
  basicstyle=\ttfamily\scriptsize,
  backgroundcolor=\color{coursepanel},frame=single,
  breaklines=true,breakatwhitespace=false,keepspaces=true,
  columns=flexible,keywordstyle=\color{courseblue}\bfseries,
  commentstyle=\color{coursegreen!80!black},
  stringstyle=\color{coursecyan},numbers=left,
  numberstyle=\tiny\color{coursemuted}
}

\newcolumntype{Y}{>{\raggedright\arraybackslash}X}
\newcolumntype{C}{>{\centering\arraybackslash}X}

\title[课程索引]{格点 QCD 核心课程全局索引}
\author[格点 QCD 核心课]{格点 QCD 核心课程讲义}
\institute{从第一性原理到梯度流核子胶子 TMD-PDF}
\date{2026 年 8 月}
\begin{document}

\begin{frame}[plain]\centering
{\LARGE\bfseries 全局索引与稳定编号\par}
\vspace{0.4cm}{\large 配合 \texttt{core\_complete.pdf} 使用\par}
\vfill 公式、算法、练习和 SymPy 证据均以 卷.单元 编号。
\end{frame}
\begin{frame}[t]{第 1 卷：数学语言、量纲与连续变化}
\small\begin{tabularx}{\linewidth}{p{0.13\linewidth}Yp{0.42\linewidth}}
\toprule 单元 & 标题 & 对象 ID\\\midrule
01.01 & 量、单位与量纲 & K/Def/Eq/Der/Fig/Alg/Ex/Sol/SYM 01.01\\
01.02 & 函数、指数与对数 & K/Def/Eq/Der/Fig/Alg/Ex/Sol/SYM 01.02\\
01.03 & 复数与相位 & K/Def/Eq/Der/Fig/Alg/Ex/Sol/SYM 01.03\\
01.04 & 导数：局部变化率 & K/Def/Eq/Der/Fig/Alg/Ex/Sol/SYM 01.04\\
01.05 & 积分：累积量与平均 & K/Def/Eq/Der/Fig/Alg/Ex/Sol/SYM 01.05\\
\bottomrule\end{tabularx}\end{frame}
\begin{frame}[t]{第 2 卷：向量、矩阵、张量与对称性}
\small\begin{tabularx}{\linewidth}{p{0.13\linewidth}Yp{0.42\linewidth}}
\toprule 单元 & 标题 & 对象 ID\\\midrule
02.01 & 向量、内积与正交 & K/Def/Eq/Der/Fig/Alg/Ex/Sol/SYM 02.01\\
02.02 & 矩阵、逆与线性方程 & K/Def/Eq/Der/Fig/Alg/Ex/Sol/SYM 02.02\\
02.03 & 本征值与谱分解 & K/Def/Eq/Der/Fig/Alg/Ex/Sol/SYM 02.03\\
02.04 & 指标、张量与缩并 & K/Def/Eq/Der/Fig/Alg/Ex/Sol/SYM 02.04\\
02.05 & 群、生成元与连续对称 & K/Def/Eq/Der/Fig/Alg/Ex/Sol/SYM 02.05\\
\bottomrule\end{tabularx}\end{frame}
\begin{frame}[t]{第 3 卷：傅里叶分析、概率与数值误差}
\small\begin{tabularx}{\linewidth}{p{0.13\linewidth}Yp{0.42\linewidth}}
\toprule 单元 & 标题 & 对象 ID\\\midrule
03.01 & 离散傅里叶变换 & K/Def/Eq/Der/Fig/Alg/Ex/Sol/SYM 03.01\\
03.02 & 高斯积分与傅里叶宽度 & K/Def/Eq/Der/Fig/Alg/Ex/Sol/SYM 03.02\\
03.03 & 概率、期望与方差 & K/Def/Eq/Der/Fig/Alg/Ex/Sol/SYM 03.03\\
03.04 & Taylor 展开与差分误差 & K/Def/Eq/Der/Fig/Alg/Ex/Sol/SYM 03.04\\
03.05 & 数值积分、收敛与误差账本 & K/Def/Eq/Der/Fig/Alg/Ex/Sol/SYM 03.05\\
\bottomrule\end{tabularx}\end{frame}
\begin{frame}[t]{第 4 卷：时空、作用量与经典场}
\small\begin{tabularx}{\linewidth}{p{0.13\linewidth}Yp{0.42\linewidth}}
\toprule 单元 & 标题 & 对象 ID\\\midrule
04.01 & Lorentz 变换与不变量 & K/Def/Eq/Der/Fig/Alg/Ex/Sol/SYM 04.01\\
04.02 & 四动量与质壳 & K/Def/Eq/Der/Fig/Alg/Ex/Sol/SYM 04.02\\
04.03 & 作用量与 Euler--Lagrange 方程 & K/Def/Eq/Der/Fig/Alg/Ex/Sol/SYM 04.03\\
04.04 & Noether 思想与守恒量 & K/Def/Eq/Der/Fig/Alg/Ex/Sol/SYM 04.04\\
04.05 & 规范势与电磁场强 & K/Def/Eq/Der/Fig/Alg/Ex/Sol/SYM 04.05\\
\bottomrule\end{tabularx}\end{frame}
\begin{frame}[t]{第 5 卷：量子态、算符与测量}
\small\begin{tabularx}{\linewidth}{p{0.13\linewidth}Yp{0.42\linewidth}}
\toprule 单元 & 标题 & 对象 ID\\\midrule
05.01 & 量子态与归一化 & K/Def/Eq/Der/Fig/Alg/Ex/Sol/SYM 05.01\\
05.02 & 算符与对易子 & K/Def/Eq/Der/Fig/Alg/Ex/Sol/SYM 05.02\\
05.03 & 本征态、两能级与时间演化 & K/Def/Eq/Der/Fig/Alg/Ex/Sol/SYM 05.03\\
05.04 & Born 规则与期望值 & K/Def/Eq/Der/Fig/Alg/Ex/Sol/SYM 05.04\\
05.05 & 不确定关系与高斯最小波包 & K/Def/Eq/Der/Fig/Alg/Ex/Sol/SYM 05.05\\
\bottomrule\end{tabularx}\end{frame}
\begin{frame}[t]{第 6 卷：从量子力学到量子场论}
\small\begin{tabularx}{\linewidth}{p{0.13\linewidth}Yp{0.42\linewidth}}
\toprule 单元 & 标题 & 对象 ID\\\midrule
06.01 & 自由标量场与 Klein--Gordon 方程 & K/Def/Eq/Der/Fig/Alg/Ex/Sol/SYM 06.01\\
06.02 & 场的正则量子化 & K/Def/Eq/Der/Fig/Alg/Ex/Sol/SYM 06.02\\
06.03 & 路径积分与生成泛函 & K/Def/Eq/Der/Fig/Alg/Ex/Sol/SYM 06.03\\
06.04 & Wick 定理与收缩 & K/Def/Eq/Der/Fig/Alg/Ex/Sol/SYM 06.04\\
06.05 & Wick 转动与欧氏谱表示 & K/Def/Eq/Der/Fig/Alg/Ex/Sol/SYM 06.05\\
\bottomrule\end{tabularx}\end{frame}
\begin{frame}[t]{第 7 卷：QCD：颜色规范理论}
\small\begin{tabularx}{\linewidth}{p{0.13\linewidth}Yp{0.42\linewidth}}
\toprule 单元 & 标题 & 对象 ID\\\midrule
07.01 & 夸克、胶子与颜色 & K/Def/Eq/Der/Fig/Alg/Ex/Sol/SYM 07.01\\
07.02 & SU(3) 生成元与李代数 & K/Def/Eq/Der/Fig/Alg/Ex/Sol/SYM 07.02\\
07.03 & 局域规范协变与协变导数 & K/Def/Eq/Der/Fig/Alg/Ex/Sol/SYM 07.03\\
07.04 & 非阿贝尔场强与 QCD 作用量 & K/Def/Eq/Der/Fig/Alg/Ex/Sol/SYM 07.04\\
07.05 & 运行耦合、渐近自由与禁闭 & K/Def/Eq/Der/Fig/Alg/Ex/Sol/SYM 07.05\\
\bottomrule\end{tabularx}\end{frame}
\begin{frame}[t]{第 8 卷：欧氏格点与规范链接}
\small\begin{tabularx}{\linewidth}{p{0.13\linewidth}Yp{0.42\linewidth}}
\toprule 单元 & 标题 & 对象 ID\\\midrule
08.01 & 有限周期格点与动量量子化 & K/Def/Eq/Der/Fig/Alg/Ex/Sol/SYM 08.01\\
08.02 & 链接变量与平行输运 & K/Def/Eq/Der/Fig/Alg/Ex/Sol/SYM 08.02\\
08.03 & Plaquette、场强与 Wilson 圈 & K/Def/Eq/Der/Fig/Alg/Ex/Sol/SYM 08.03\\
08.04 & Wilson 规范作用量与连续极限 & K/Def/Eq/Der/Fig/Alg/Ex/Sol/SYM 08.04\\
08.05 & 尺度、有限体积与 Symanzik 展开 & K/Def/Eq/Der/Fig/Alg/Ex/Sol/SYM 08.05\\
\bottomrule\end{tabularx}\end{frame}
\begin{frame}[t]{第 9 卷：格点费米子、传播子与涂抹}
\small\begin{tabularx}{\linewidth}{p{0.13\linewidth}Yp{0.42\linewidth}}
\toprule 单元 & 标题 & 对象 ID\\\midrule
09.01 & Dirac 旋量与 \ensuremath{\gamma} 矩阵 & K/Def/Eq/Der/Fig/Alg/Ex/Sol/SYM 09.01\\
09.02 & 朴素离散与费米子倍增 & K/Def/Eq/Der/Fig/Alg/Ex/Sol/SYM 09.02\\
09.03 & Wilson 与 Clover 改进 & K/Def/Eq/Der/Fig/Alg/Ex/Sol/SYM 09.03\\
09.04 & 传播子与迭代线性求解 & K/Def/Eq/Der/Fig/Alg/Ex/Sol/SYM 09.04\\
09.05 & 夸克场涂抹与高动量源 & K/Def/Eq/Der/Fig/Alg/Ex/Sol/SYM 09.05\\
\bottomrule\end{tabularx}\end{frame}
\begin{frame}[t]{第 10 卷：Monte Carlo、HMC 与规范系综}
\small\begin{tabularx}{\linewidth}{p{0.13\linewidth}Yp{0.42\linewidth}}
\toprule 单元 & 标题 & 对象 ID\\\midrule
10.01 & 欧氏权重与重要性抽样 & K/Def/Eq/Der/Fig/Alg/Ex/Sol/SYM 10.01\\
10.02 & Metropolis--Hastings 与详细平衡 & K/Def/Eq/Der/Fig/Alg/Ex/Sol/SYM 10.02\\
10.03 & 热化、自相关与有效样本数 & K/Def/Eq/Der/Fig/Alg/Ex/Sol/SYM 10.03\\
10.04 & Hybrid Monte Carlo & K/Def/Eq/Der/Fig/Alg/Ex/Sol/SYM 10.04\\
10.05 & 系综元数据、尺度与质量门 & K/Def/Eq/Der/Fig/Alg/Ex/Sol/SYM 10.05\\
\bottomrule\end{tabularx}\end{frame}
\begin{frame}[t]{第 11 卷：核子插值算符与二、三点关联函数}
\small\begin{tabularx}{\linewidth}{p{0.13\linewidth}Yp{0.42\linewidth}}
\toprule 单元 & 标题 & 对象 ID\\\midrule
11.01 & 质子插值算符与量子数 & K/Def/Eq/Der/Fig/Alg/Ex/Sol/SYM 11.01\\
11.02 & 核子二点函数与谱分解 & K/Def/Eq/Der/Fig/Alg/Ex/Sol/SYM 11.02\\
11.03 & 有效质量与相关拟合 & K/Def/Eq/Der/Fig/Alg/Ex/Sol/SYM 11.03\\
11.04 & 三点函数、比值与激发态 & K/Def/Eq/Der/Fig/Alg/Ex/Sol/SYM 11.04\\
11.05 & 断连图与真空扣除 & K/Def/Eq/Der/Fig/Alg/Ex/Sol/SYM 11.05\\
\bottomrule\end{tabularx}\end{frame}
\begin{frame}[t]{第 12 卷：重采样、协方差与多态拟合}
\small\begin{tabularx}{\linewidth}{p{0.13\linewidth}Yp{0.42\linewidth}}
\toprule 单元 & 标题 & 对象 ID\\\midrule
12.01 & 样本均值、协方差与相关性 & K/Def/Eq/Der/Fig/Alg/Ex/Sol/SYM 12.01\\
12.02 & Jackknife 与非线性观测量 & K/Def/Eq/Der/Fig/Alg/Ex/Sol/SYM 12.02\\
12.03 & Bootstrap、分位数与非高斯分布 & K/Def/Eq/Der/Fig/Alg/Ex/Sol/SYM 12.03\\
12.04 & 相关 \ensuremath{\chi}\ensuremath{^{2}} 与协方差正则化 & K/Def/Eq/Der/Fig/Alg/Ex/Sol/SYM 12.04\\
12.05 & 多态联合拟合与模型系统学 & K/Def/Eq/Der/Fig/Alg/Ex/Sol/SYM 12.05\\
\bottomrule\end{tabularx}\end{frame}
\begin{frame}[t]{第 13 卷：深度非弹性散射与部分子分布}
\small\begin{tabularx}{\linewidth}{p{0.13\linewidth}Yp{0.42\linewidth}}
\toprule 单元 & 标题 & 对象 ID\\\midrule
13.01 & DIS 运动学与 Bjorken x & K/Def/Eq/Der/Fig/Alg/Ex/Sol/SYM 13.01\\
13.02 & 光锥 PDF 的算符定义 & K/Def/Eq/Der/Fig/Alg/Ex/Sol/SYM 13.02\\
13.03 & OPE、Mellin 矩与局域算符 & K/Def/Eq/Der/Fig/Alg/Ex/Sol/SYM 13.03\\
13.04 & DGLAP 演化与尺度依赖 & K/Def/Eq/Der/Fig/Alg/Ex/Sol/SYM 13.04\\
13.05 & 从 collinear PDF 到 TMD & K/Def/Eq/Der/Fig/Alg/Ex/Sol/SYM 13.05\\
\bottomrule\end{tabularx}\end{frame}
\begin{frame}[t]{第 14 卷：LaMET、quasi 分布与 Collins--Soper 演化}
\small\begin{tabularx}{\linewidth}{p{0.13\linewidth}Yp{0.42\linewidth}}
\toprule 单元 & 标题 & 对象 ID\\\midrule
14.01 & 欧氏等时相关与 quasi-PDF & K/Def/Eq/Der/Fig/Alg/Ex/Sol/SYM 14.01\\
14.02 & LaMET 因子化与幂修正 & K/Def/Eq/Der/Fig/Alg/Ex/Sol/SYM 14.02\\
14.03 & 匹配卷积与胶子—夸克混合 & K/Def/Eq/Der/Fig/Alg/Ex/Sol/SYM 14.03\\
14.04 & 伪分布与 Ioffe-time 分布 & K/Def/Eq/Der/Fig/Alg/Ex/Sol/SYM 14.04\\
14.05 & Collins--Soper 核与多动量比值 & K/Def/Eq/Der/Fig/Alg/Ex/Sol/SYM 14.05\\
\bottomrule\end{tabularx}\end{frame}
\begin{frame}[t]{第 15 卷：胶子 TMD 算符、投影与 staple 几何}
\small\begin{tabularx}{\linewidth}{p{0.13\linewidth}Yp{0.42\linewidth}}
\toprule 单元 & 标题 & 对象 ID\\\midrule
15.01 & Clover 场强与梯度流后胶子场 & K/Def/Eq/Der/Fig/Alg/Ex/Sol/SYM 15.01\\
15.02 & 胶子双场强算符与规范不变性 & K/Def/Eq/Der/Fig/Alg/Ex/Sol/SYM 15.02\\
15.03 & 非零 bT 的有限 staple 路径 & K/Def/Eq/Der/Fig/Alg/Ex/Sol/SYM 15.03\\
15.04 & 非极化与极化张量投影 & K/Def/Eq/Der/Fig/Alg/Ex/Sol/SYM 15.04\\
15.05 & 算符基底、混合与接口契约 & K/Def/Eq/Der/Fig/Alg/Ex/Sol/SYM 15.05\\
\bottomrule\end{tabularx}\end{frame}
\begin{frame}[t]{第 16 卷：梯度流与流时重整化方案}
\small\begin{tabularx}{\linewidth}{p{0.13\linewidth}Yp{0.42\linewidth}}
\toprule 单元 & 标题 & 对象 ID\\\midrule
16.01 & 热方程、热核与平滑半径 & K/Def/Eq/Der/Fig/Alg/Ex/Sol/SYM 16.01\\
16.02 & Yang--Mills/\allowbreak{}Wilson 流方程 & K/Def/Eq/Der/Fig/Alg/Ex/Sol/SYM 16.02\\
16.03 & 流时窗口与尺度分离 & K/Def/Eq/Der/Fig/Alg/Ex/Sol/SYM 16.03\\
16.04 & 有限流时复合算符与重整化边界 & K/Def/Eq/Der/Fig/Alg/Ex/Sol/SYM 16.04\\
16.05 & 局域小流时展开与非局域路径边界 & K/Def/Eq/Der/Fig/Alg/Ex/Sol/SYM 16.05\\
\bottomrule\end{tabularx}\end{frame}
\begin{frame}[t]{第 17 卷：核子中 flowed 胶子 TMD 的格点测量}
\small\begin{tabularx}{\linewidth}{p{0.13\linewidth}Yp{0.42\linewidth}}
\toprule 单元 & 标题 & 对象 ID\\\midrule
17.01 & 从原始链接到 flowed 场强 & K/Def/Eq/Der/Fig/Alg/Ex/Sol/SYM 17.01\\
17.02 & 批量生成 z、bT、\ensuremath{\ell} 与投影 & K/Def/Eq/Der/Fig/Alg/Ex/Sol/SYM 17.02\\
17.03 & boost 核子二点函数与动量涂抹 & K/Def/Eq/Der/Fig/Alg/Ex/Sol/SYM 17.03\\
17.04 & 构型级断连三点与真空扣除 & K/Def/Eq/Der/Fig/Alg/Ex/Sol/SYM 17.04\\
17.05 & 平台、summation 与两态联合分析 & K/Def/Eq/Der/Fig/Alg/Ex/Sol/SYM 17.05\\
\bottomrule\end{tabularx}\end{frame}
\begin{frame}[t]{第 18 卷：soft、ZR、hybrid 重整化与匹配}
\small\begin{tabularx}{\linewidth}{p{0.13\linewidth}Yp{0.42\linewidth}}
\toprule 单元 & 标题 & 对象 ID\\\midrule
18.01 & Wilson 线自能、线性发散与流时 & K/Def/Eq/Der/Fig/Alg/Ex/Sol/SYM 18.01\\
18.02 & 方案化 quasi-soft、rapidity 扣除与标准 TMD 转换 & K/Def/Eq/Der/Fig/Alg/Ex/Sol/SYM 18.02\\
18.03 & ZR 比值重整化 & K/Def/Eq/Der/Fig/Alg/Ex/Sol/SYM 18.03\\
18.04 & 直线算符的 hybrid 重整化与 staple 推广边界 & K/Def/Eq/Der/Fig/Alg/Ex/Sol/SYM 18.04\\
18.05 & 从重整化坐标矩阵元到物理胶子 TMD & K/Def/Eq/Der/Fig/Alg/Ex/Sol/SYM 18.05\\
\bottomrule\end{tabularx}\end{frame}
\begin{frame}[t]{第 19 卷：连续、大动量、零流时与无限 staple 联合极限}
\small\begin{tabularx}{\linewidth}{p{0.13\linewidth}Yp{0.42\linewidth}}
\toprule 单元 & 标题 & 对象 ID\\\midrule
19.01 & 连续极限 a\ensuremath{\rightarrow}0 & K/Def/Eq/Der/Fig/Alg/Ex/Sol/SYM 19.01\\
19.02 & 大动量极限 Pz\ensuremath{\rightarrow}\ensuremath{\infty} & K/Def/Eq/Der/Fig/Alg/Ex/Sol/SYM 19.02\\
19.03 & 零流时/\allowbreak{}方案转换：局域模型与 staple 停止门 & K/Def/Eq/Der/Fig/Alg/Ex/Sol/SYM 19.03\\
19.04 & 有限 staple 长度与边界极限 & K/Def/Eq/Der/Fig/Alg/Ex/Sol/SYM 19.04\\
19.05 & 联合拟合、可辨识性与总误差预算 & K/Def/Eq/Der/Fig/Alg/Ex/Sol/SYM 19.05\\
\bottomrule\end{tabularx}\end{frame}
\begin{frame}[t]{第 20 卷：自主实现：从数据契约到可发布结果}
\small\begin{tabularx}{\linewidth}{p{0.13\linewidth}Yp{0.42\linewidth}}
\toprule 单元 & 标题 & 对象 ID\\\midrule
20.01 & 物理对象到数组与元数据契约 & K/Def/Eq/Der/Fig/Alg/Ex/Sol/SYM 20.01\\
20.02 & 六阶段流水线与断点续跑 & K/Def/Eq/Der/Fig/Alg/Ex/Sol/SYM 20.02\\
20.03 & 分层测试：代数、几何、统计与物理闭合 & K/Def/Eq/Der/Fig/Alg/Ex/Sol/SYM 20.03\\
20.04 & 并行、I/\allowbreak{}O、性能与可复现运行 & K/Def/Eq/Der/Fig/Alg/Ex/Sol/SYM 20.04\\
20.05 & 结课项目：生产判据与发布审计 & K/Def/Eq/Der/Fig/Alg/Ex/Sol/SYM 20.05\\
\bottomrule\end{tabularx}\end{frame}
\begin{frame}[t]{第 21 卷：有限群、表示与格点对称性}
\small\begin{tabularx}{\linewidth}{p{0.13\linewidth}Yp{0.42\linewidth}}
\toprule 单元 & 标题 & 对象 ID\\\midrule
21.01 & 群作用、子群、陪集与 Lagrange 定理 & K/Def/Eq/Der/Fig/Alg/Ex/Sol/SYM 21.01\\
21.02 & 表示、不可约表示与 Schur 引理三情形 & K/Def/Eq/Der/Fig/Alg/Ex/Sol/SYM 21.02\\
21.03 & 特征标正交与表示分解 & K/Def/Eq/Der/Fig/Alg/Ex/Sol/SYM 21.03\\
21.04 & 立方旋转群 O、O\_h、宇称与自旋下沉 & K/Def/Eq/Der/Fig/Alg/Ex/Sol/SYM 21.04\\
21.05 & 有限群投影算符 & K/Def/Eq/Der/Fig/Alg/Ex/Sol/SYM 21.05\\
\bottomrule\end{tabularx}\end{frame}
\begin{frame}[t]{第 22 卷：李群、李代数与 SU(3) 表示}
\small\begin{tabularx}{\linewidth}{p{0.13\linewidth}Yp{0.42\linewidth}}
\toprule 单元 & 标题 & 对象 ID\\\midrule
22.01 & 微分几何桥梁：流形、切空间与指数映射 & K/Def/Eq/Der/Fig/Alg/Ex/Sol/SYM 22.01\\
22.02 & 厄米生成元、结构常数、Jacobi 恒等式与 BCH & K/Def/Eq/Der/Fig/Alg/Ex/Sol/SYM 22.02\\
22.03 & SU(2) 角动量表示与 Casimir & K/Def/Eq/Der/Fig/Alg/Ex/Sol/SYM 22.03\\
22.04 & SU(3) 根、权、Dynkin 标号、维数与 Casimir & K/Def/Eq/Der/Fig/Alg/Ex/Sol/SYM 22.04\\
22.05 & 测度桥梁、Haar 积分与矩阵元正交 & K/Def/Eq/Der/Fig/Alg/Ex/Sol/SYM 22.05\\
\bottomrule\end{tabularx}\end{frame}
\begin{frame}[t]{第 23 卷：粒子物理、标准模型与强子量子数}
\small\begin{tabularx}{\linewidth}{p{0.13\linewidth}Yp{0.42\linewidth}}
\toprule 单元 & 标题 & 对象 ID\\\midrule
23.01 & 一代标准模型场表、手征性与电荷 & K/Def/Eq/Der/Fig/Alg/Ex/Sol/SYM 23.01\\
23.02 & 夸克模型、flavor 多重态与强子分类 & K/Def/Eq/Der/Fig/Alg/Ex/Sol/SYM 23.02\\
23.03 & 角动量耦合、宇称、荷共轭与 JPC & K/Def/Eq/Der/Fig/Alg/Ex/Sol/SYM 23.03\\
23.04 & 两体相空间、散射截面与衰变宽度 & K/Def/Eq/Der/Fig/Alg/Ex/Sol/SYM 23.04\\
23.05 & 显式卷积、硬散射因子化与部分子模型 & K/Def/Eq/Der/Fig/Alg/Ex/Sol/SYM 23.05\\
\bottomrule\end{tabularx}\end{frame}
\begin{frame}[t]{第 24 卷：量子场论深化：传播子、微扰论与 LSZ}
\small\begin{tabularx}{\linewidth}{p{0.13\linewidth}Yp{0.42\linewidth}}
\toprule 单元 & 标题 & 对象 ID\\\midrule
24.01 & 正则量子化、模展开与真空 & K/Def/Eq/Der/Fig/Alg/Ex/Sol/SYM 24.01\\
24.02 & 复分析与分布桥梁：传播子、留数和 i\ensuremath{\epsilon} & K/Def/Eq/Der/Fig/Alg/Ex/Sol/SYM 24.02\\
24.03 & 相互作用图、生成泛函与 Feynman 规则 & K/Def/Eq/Der/Fig/Alg/Ex/Sol/SYM 24.03\\
24.04 & 固定归一化的 Minkowski LSZ 与 Euclidean 边界 & K/Def/Eq/Der/Fig/Alg/Ex/Sol/SYM 24.04\\
24.05 & \ensuremath{\lambda}\ensuremath{\phi}\ensuremath{^{4}} 一圈：正则化、极点、反项与 beta & K/Def/Eq/Der/Fig/Alg/Ex/Sol/SYM 24.05\\
\bottomrule\end{tabularx}\end{frame}
\begin{frame}[t]{第 25 卷：规范场论深化：规范固定、鬼场、BRST 与反常}
\small\begin{tabularx}{\linewidth}{p{0.13\linewidth}Yp{0.42\linewidth}}
\toprule 单元 & 标题 & 对象 ID\\\midrule
25.01 & Hamilton 桥梁：相空间、约束与两个物理偏振 & K/Def/Eq/Der/Fig/Alg/Ex/Sol/SYM 25.01\\
25.02 & Grassmann 桥梁与 Faddeev--Popov 行列式 & K/Def/Eq/Der/Fig/Alg/Ex/Sol/SYM 25.02\\
25.03 & 完整 BRST、分次 Leibniz 与规范固定作用量 & K/Def/Eq/Der/Fig/Alg/Ex/Sol/SYM 25.03\\
25.04 & Z1、Z2、Ward 与 Slavnov--Taylor 恒等式 & K/Def/Eq/Der/Fig/Alg/Ex/Sol/SYM 25.04\\
25.05 & 轴反常统一约定与一代标准模型反常抵消 & K/Def/Eq/Der/Fig/Alg/Ex/Sol/SYM 25.05\\
\bottomrule\end{tabularx}\end{frame}
\begin{frame}[t]{第 26 卷：非微扰场论：对称破缺、有效理论、拓扑与禁闭}
\small\begin{tabularx}{\linewidth}{p{0.13\linewidth}Yp{0.42\linewidth}}
\toprule 单元 & 标题 & 对象 ID\\\midrule
26.01 & 自发对称破缺与 Goldstone 定理 & K/Def/Eq/Der/Fig/Alg/Ex/Sol/SYM 26.01\\
26.02 & Higgs 机制与规范玻色子质量 & K/Def/Eq/Der/Fig/Alg/Ex/Sol/SYM 26.02\\
26.03 & 有效场论、OPE、matching 与系统幂计数 & K/Def/Eq/Der/Fig/Alg/Ex/Sol/SYM 26.03\\
26.04 & 拓扑桥梁：S\ensuremath{^{3}}\ensuremath{\infty}\ensuremath{\rightarrow}SU(N)、\ensuremath{\pi}3=Z 与 instanton & K/Def/Eq/Der/Fig/Alg/Ex/Sol/SYM 26.04\\
26.05 & Wilson 圈、Polyakov loop 中心对称与禁闭边界 & K/Def/Eq/Der/Fig/Alg/Ex/Sol/SYM 26.05\\
\bottomrule\end{tabularx}\end{frame}
\begin{frame}[t]{第 27 卷：格点谱学：从关联函数到强子能谱}
\small\begin{tabularx}{\linewidth}{p{0.13\linewidth}Yp{0.42\linewidth}}
\toprule 单元 & 标题 & 对象 ID\\\midrule
27.01 & 转移矩阵与二点函数谱分解 & K/Def/Eq/Der/Fig/Alg/Ex/Sol/SYM 27.01\\
27.02 & 有效质量、周期边界与多态拟合 & K/Def/Eq/Der/Fig/Alg/Ex/Sol/SYM 27.02\\
27.03 & 相关矩阵与广义本征值问题 & K/Def/Eq/Der/Fig/Alg/Ex/Sol/SYM 27.03\\
27.04 & LapH distillation、perambulator 与算符因子化 & K/Def/Eq/Der/Fig/Alg/Ex/Sol/SYM 27.04\\
27.05 & 立方群下沉、算符重叠与连续自旋指认 & K/Def/Eq/Der/Fig/Alg/Ex/Sol/SYM 27.05\\
\bottomrule\end{tabularx}\end{frame}
\begin{frame}[t]{第 28 卷：多强子谱、Lüscher 方法与共振极点}
\small\begin{tabularx}{\linewidth}{p{0.13\linewidth}Yp{0.42\linewidth}}
\toprule 单元 & 标题 & 对象 ID\\\midrule
28.01 & 两粒子运动学与非相互作用能级 & K/Def/Eq/Der/Fig/Alg/Ex/Sol/SYM 28.01\\
28.02 & 弹性 Lüscher 量化条件与相移 & K/Def/Eq/Der/Fig/Alg/Ex/Sol/SYM 28.02\\
28.03 & 有效程展开、散射长度与束缚态 & K/Def/Eq/Der/Fig/Alg/Ex/Sol/SYM 28.03\\
28.04 & moving frame、little group 与分波混合 & K/Def/Eq/Der/Fig/Alg/Ex/Sol/SYM 28.04\\
28.05 & 耦合道振幅、解析延拓与共振极点 & K/Def/Eq/Der/Fig/Alg/Ex/Sol/SYM 28.05\\
\bottomrule\end{tabularx}\end{frame}
\begin{frame}[t]{第 29 卷：格点有限体积系统学}
\small\begin{tabularx}{\linewidth}{p{0.13\linewidth}Yp{0.42\linewidth}}
\toprule 单元 & 标题 & 对象 ID\\\midrule
29.01 & Poisson 求和与有限盒修正的来源 & K/Def/Eq/Der/Fig/Alg/Ex/Sol/SYM 29.01\\
29.02 & 质量隙、m\ensuremath{\pi}L 与单强子指数修正 & K/Def/Eq/Der/Fig/Alg/Ex/Sol/SYM 29.02\\
29.03 & 两粒子在壳传播与幂律体积效应 & K/Def/Eq/Der/Fig/Alg/Ex/Sol/SYM 29.03\\
29.04 & 扭曲边界条件与连续动量调谐 & K/Def/Eq/Der/Fig/Alg/Ex/Sol/SYM 29.04\\
29.05 & QED 有限体积、零模与方案依赖 & K/Def/Eq/Der/Fig/Alg/Ex/Sol/SYM 29.05\\
\bottomrule\end{tabularx}\end{frame}
\begin{frame}[t]{第 30 卷：有限温度格点 QCD 与热介质}
\small\begin{tabularx}{\linewidth}{p{0.13\linewidth}Yp{0.42\linewidth}}
\toprule 单元 & 标题 & 对象 ID\\\midrule
30.01 & 热迹、Euclidean 时间圆与温度设定 & K/Def/Eq/Der/Fig/Alg/Ex/Sol/SYM 30.01\\
30.02 & Polyakov loop、中心对称与自由能 & K/Def/Eq/Der/Fig/Alg/Ex/Sol/SYM 30.02\\
30.03 & 手征凝聚、susceptibility 与 crossover & K/Def/Eq/Der/Fig/Alg/Ex/Sol/SYM 30.03\\
30.04 & 状态方程、trace anomaly 与积分法 & K/Def/Eq/Der/Fig/Alg/Ex/Sol/SYM 30.04\\
30.05 & 热谱函数、输运系数与病态反演 & K/Def/Eq/Der/Fig/Alg/Ex/Sol/SYM 30.05\\
\bottomrule\end{tabularx}\end{frame}
\begin{frame}[t]{第 31 卷：格点费米子方案全景}
\small\begin{tabularx}{\linewidth}{p{0.13\linewidth}Yp{0.42\linewidth}}
\toprule 单元 & 标题 & 对象 ID\\\midrule
31.01 & staggered 费米子、taste 与 rooting & K/Def/Eq/Der/Fig/Alg/Ex/Sol/SYM 31.01\\
31.02 & domain-wall 费米子与剩余质量 & K/Def/Eq/Der/Fig/Alg/Ex/Sol/SYM 31.02\\
31.03 & overlap 算子与 Ginsparg--Wilson 关系 & K/Def/Eq/Der/Fig/Alg/Ex/Sol/SYM 31.03\\
31.04 & twisted-mass Wilson 与最大扭转自动改进 & K/Def/Eq/Der/Fig/Alg/Ex/Sol/SYM 31.04\\
31.05 & mixed-action、partial quenching 与统一连续极限 & K/Def/Eq/Der/Fig/Alg/Ex/Sol/SYM 31.05\\
\bottomrule\end{tabularx}\end{frame}
\begin{frame}[t]{第 32 卷：改进规范作用量、各向异性、拓扑与尺度设定}
\small\begin{tabularx}{\linewidth}{p{0.13\linewidth}Yp{0.42\linewidth}}
\toprule 单元 & 标题 & 对象 ID\\\midrule
32.01 & Symanzik 改进规范作用量 & K/Def/Eq/Der/Fig/Alg/Ex/Sol/SYM 32.01\\
32.02 & 各向异性格点与色散调谐 & K/Def/Eq/Der/Fig/Alg/Ex/Sol/SYM 32.02\\
32.03 & 拓扑冻结、自相关与算法临界变慢 & K/Def/Eq/Der/Fig/Alg/Ex/Sol/SYM 32.03\\
32.04 & 开放时间边界与边界污染 & K/Def/Eq/Der/Fig/Alg/Ex/Sol/SYM 32.04\\
32.05 & 尺度设定：r0、t0、w0 与误差传播 & K/Def/Eq/Der/Fig/Alg/Ex/Sol/SYM 32.05\\
\bottomrule\end{tabularx}\end{frame}
\begin{frame}[t]{第 33 卷：局域算符重整化方案与非微扰演化}
\small\begin{tabularx}{\linewidth}{p{0.13\linewidth}Yp{0.42\linewidth}}
\toprule 单元 & 标题 & 对象 ID\\\midrule
33.01 & 重整化矩阵、算符混合与 MSbar & K/Def/Eq/Der/Fig/Alg/Ex/Sol/SYM 33.01\\
33.02 & RI/\allowbreak{}MOM：固定规范下的非微扰重整化 & K/Def/Eq/Der/Fig/Alg/Ex/Sol/SYM 33.02\\
33.03 & RI/\allowbreak{}SMOM：非例外动量与 Ward 恒等式 & K/Def/Eq/Der/Fig/Alg/Ex/Sol/SYM 33.03\\
33.04 & Schrödinger functional 与算符 step scaling & K/Def/Eq/Der/Fig/Alg/Ex/Sol/SYM 33.04\\
33.05 & gradient-flow coupling 与有限体积 beta function & K/Def/Eq/Der/Fig/Alg/Ex/Sol/SYM 33.05\\
\bottomrule\end{tabularx}\end{frame}
\begin{frame}[t]{第 34 卷：非局域算符与 TMD 重整化方案}
\small\begin{tabularx}{\linewidth}{p{0.13\linewidth}Yp{0.42\linewidth}}
\toprule 单元 & 标题 & 对象 ID\\\midrule
34.01 & Wilson-line 线性发散与乘法结构 & K/Def/Eq/Der/Fig/Alg/Ex/Sol/SYM 34.01\\
34.02 & 辅助场方法：Wilson line 化为静态传播子 & K/Def/Eq/Der/Fig/Alg/Ex/Sol/SYM 34.02\\
34.03 & ratio 与自重整化：从同源数据消去 UV 因子 & K/Def/Eq/Der/Fig/Alg/Ex/Sol/SYM 34.03\\
34.04 & hybrid 方案、切换尺度与连续拼接 & K/Def/Eq/Der/Fig/Alg/Ex/Sol/SYM 34.04\\
34.05 & 方案标记的 quasi-TMD、soft、CS kernel 与 flow 转换 & K/Def/Eq/Der/Fig/Alg/Ex/Sol/SYM 34.05\\
\bottomrule\end{tabularx}\end{frame}
\begin{frame}[t]{第 35 卷：研究级联合设计、可复现审计与最终资格考核}
\small\begin{tabularx}{\linewidth}{p{0.13\linewidth}Yp{0.42\linewidth}}
\toprule 单元 & 标题 & 对象 ID\\\midrule
35.01 & 从 \ensuremath{f_1^{g[-,-]}} estimand 到四级证据状态、DAG 与停止门 & K/Def/Eq/Der/Fig/Alg/Ex/Sol/SYM 35.01\\
35.02 & flow window 与 a、\ensuremath{\tau}、L、Pz、ell 的分阶段层级极限 & K/Def/Eq/Der/Fig/Alg/Ex/Sol/SYM 35.02\\
35.03 & 有限 Ioffe-time 数据的 Fourier 逆问题 & K/Def/Eq/Der/Fig/Alg/Ex/Sol/SYM 35.03\\
35.04 & 可复现管线、数据谱系与审计 & K/Def/Eq/Der/Fig/Alg/Ex/Sol/SYM 35.04\\
35.05 & 最终资格考核：按四级证据独立实现 \ensuremath{f_1^{g[-,-]}} 链 & K/Def/Eq/Der/Fig/Alg/Ex/Sol/SYM 35.05\\
\bottomrule\end{tabularx}\end{frame}
\end{document}
